\subsection{Statements and Well-Formed Language}

\begin{definition} \label{def:mathematical-statements-and-truth-values} \hlt{Mathematical Statements and Truth Values}\\
The \hlt{truth values} used here are \hlt{true} and \hlt{false}.\\
A \hlt{mathematical statement} is a declarative sentence having exactly one truth value.
\end{definition}

\begin{principle} \label{princ:well-formed-statements} \hlt{Well-Formed Statements}\\
A mathematical expression may serve as a statement only when all the following conditions hold.
\begin{enumerate}[label=(\roman*)]
\setlength{\itemsep}{0pt}
\item \hlt{Symbols}: All symbols have declared syntactic roles.
\item \hlt{Meaning}: All constituent expressions have defined meanings.
\item \hlt{Scope}: All variable occurrences have unambiguous scope.
\end{enumerate}
Failure of any condition makes the expression ill formed rather than false; no truth value is assigned.
\end{principle}
